%%%%%%%%%%%%%%%%%%%%%%%%%%%%%%%%%%%%%%%%%%%%%%%%%%%%%%%%%%%%%%%%%%%%%%%%%%%
% Función: ejemplos de figuras y tablas.
%
% Uso: modifica este fichero para adaptar el contenido a tu trabajo.
%%%%%%%%%%%%%%%%%%%%%%%%%%%%%%%%%%%%%%%%%%%%%%%%%%%%%%%%%%%%%%%%%%%%%%%%%%%

\section{Figuras y tablas}
\label{sec:figuras-tablas}


\subsection{Figuras de uso habitual}
\label{sec:figuras-habituales}

Esta sección presenta ejemplos de figuras y tablas que pueden adaptarse a la mayoría de los documentos. Los casos que requieren una composición más elaborada se reúnen en el capítulo~\ref{cha:ejemplos-avanzados}.

La figura~\ref{fig:figura-convencional} muestra la inclusión habitual de una imagen con \texttt{figure}, \texttt{\textbackslash{}includegraphics}, un pie descriptivo y una etiqueta para citarla desde el texto.

\begin{figure}[h]
  \centering
  \includegraphics[width=0.55\textwidth]{Figure1}
  \caption{Ejemplo de figura convencional.}
  \label{fig:figura-convencional}
\end{figure}

\subsection{Tabla convencional}
\label{sec:tabla-convencional}

La tabla~\ref{tab:table1} muestra una estructura convencional con título, etiqueta y referencia cruzada.

\begin{table}
  % increase table row spacing, adjust to taste
  \renewcommand{\arraystretch}{1.3}
  \caption{Comparativa.}
  \label{tab:table1}
  \begin{center}
    % Some packages, such as MDW tools, offer better commands for making tables than the plain LaTeX2e tabular which is used here.
    \begin{tabular}{|c|c|c|}
      \hline
      Method & Training Time & Man-Work (\%)\\
      \hline
      Propagation model & $<$ 30 sec & 5\\
      \hline
      Manual & 9 h 30 min & 24\\
      \hline
      Automatic & 2 h & 10 8\\
      \hline
    \end{tabular}
  \end{center}
\end{table}

\subsection{Subfiguras}
\label{sec:subfiguras}

El entorno \texttt{subfigure} permite reunir varias imágenes dentro de una misma figura. La figura~\ref{fig:fig3} contiene dos subfiguras, que también pueden citarse individualmente; por ejemplo, \ref{fig:fig3}.\subref{fig:fig3b}.

% For this to work you need to (in preamble.tex):
% - remove \usepackage{subfig}
% - add \usepackage{caption}
% - add \usepackage{subcaption}
\begin{figure}
  \centering
  \begin{subfigure}[b]{0.3\textwidth}
    \includegraphics[width=\textwidth]{Figure2}
    \caption{Mean Entropy.}
    \label{fig:fig3a}
  \end{subfigure}%
  ~ %add desired spacing between images, e. g. ~, \quad, \qquad etc.
  % (or a blank line to force the subfigure onto a new line)
  \begin{subfigure}[b]{0.3\textwidth}
    \includegraphics[width=\textwidth]{Figure3}
    \caption{Error Percentage.}
    \label{fig:fig3b}
  \end{subfigure}
  \caption{Optimal trames number in the training data set.}
  \label{fig:fig3}
\end{figure}

Los ejemplos avanzados de figuras y tablas, incluidas las tablas multipágina y apaisadas, las figuras rodeadas por texto y las composiciones con numerosas subfiguras, se encuentran en el capítulo~\ref{cha:ejemplos-avanzados}.


%%% Local Variables:
%%% TeX-master: "../book"
%%% End:
